\documentclass[aspectratio=169,t,13pt,usenames,dvipsnames]{beamer}
\usetheme[lily]{PaloAlto}
\usecolortheme{dolphin}
\usepackage{mathtools}
\usepackage{tikz}
\usepackage{pgfplots}
\usepackage{tikz-3dplot}
\usepackage{adjustbox}
\usetikzlibrary{quantikz}
\usepackage{xcolor}
\pgfplotsset{compat=newest}
\logo{%
\stackbox[c][c]{\includegraphics[height=1.15cm]{logo.pdf}\\[-3pt]{\textcopyright~2021}\\{Ron K. Cytron}}}

% \logo{\includegraphics[height=2.5cm]{logo.pdf}}
\def\ColorR#1{\textcolor{BrickRed}{#1}}
\def\ColorG#1{\textcolor{OliveGreen}{#1}}
\def\ColorB#1{\textcolor{NavyBlue}{#1}}
\definecolor{links}{HTML}{2E8B57}
\hypersetup{colorlinks,linkcolor=,urlcolor=links}
\def\LinkArrow#1{%
\mbox{$\,$\href{#1}{\mbox{\vrule width 0pt height 0ex depth 0pt$\mapsto$}}}$\,$}
\def\True{\mbox{\texttt{True}}}
\def\False{\mbox{\texttt{False}}}
\def\Zero{\mbox{$0$}}
\def\One{\mbox{$1$}}
\def\Not#1{%
\ensuremath{\Overline{#1}}}
\def\Xor#1#2{\ensuremath{#1\oplus #2}}
\def\Nand#1#2{%
\ensuremath{\mbox{Nand}(#1,#2)}}
\def\And#1#2{%
\ensuremath{#1 \wedge #2}}
\def\Or#1#2{%
\ensuremath{#1 \vee #2}}
\def\Conj#1{%
\ensuremath{#1^{\star}}}
\def\Star#1{%
\Conj{#1}}
\def\EmptyString{\ensuremath{\lambda}}
\def\Mag#1{%
\ensuremath{|\,#1\,|}}
\def\Prob#1{%
\ensuremath{\Mag{#1}^{2}}}
\def\CNOT#1#2{%
\ensuremath{\mbox{CNOT}(#1,#2)}}
\def\CNOTMatrix{%
\begin{pmatrix*}[r]
1 & 0 & 0 & 0 \\
0 & 1 & 0 & 0 \\
0 & 0 & 0 & 1 \\
0 & 0 & 1 & 0\end{pmatrix*}%
}
\def\CCNOT#1#2#3{%
\ensuremath{\mbox{CNOT}(#1,#2,#3)}}
\long\def\ToDo#1{\smash{\textcolor{YellowOrange}{#1}}}
\def\ApplyGate#1#2{%
\ensuremath{#1\,#2}}
\def\Inverse#1{\ensuremath{{#1}^{-1}}}
\def\Quote#1{``#1''}
\def\NamedGate#1{%
\mbox{\bfseries{\textrm{#1}}}}
\def\Identity{\NamedGate{I}}
\def\Hadamard{\NamedGate{H}}
\def\PauliX{\NamedGate{X}}
\def\PauliY{\NamedGate{Y}}
\def\PauliZ{\NamedGate{Z}}
\def\SQBG#1#2#3#4#5{%
\ensuremath{%
#1 \begin{pmatrix*}[r] #2 & #3 \\ #4 & #5\end{pmatrix*}}}
\def\HMatrix{%
\SQBG{\RootTwo{}}{1}{1}{1}{-1}}
\def\XMatrix{%
\SQBG{\relax}{0}{1}{1}{0}}
\def\ZMatrix{%
\SQBG{\relax}{1}{0}{0}{-1}}
\def\YMatrix{%
\SQBG{\relax}{0}{-i}{i}{0}}
\def\Co#1{\mathbb{C}}
\def\Implies#1#2{\ensuremath{#1 \rightarrow\ #2}}
\def\Domain#1{\ensuremath{\mathcal{#1}}}
\def\CompClass#1{\Domain{#1}}
\def\SpecialX#1{\ensuremath{#1^{\star}}}
\def\Forall#1#2{\ensuremath{\forall\ #1, #2}}
\def\Degrees#1{\ensuremath{#1^{\circ}}}
\def\RootTwo{\ensuremath{\frac{1}{\sqrt{2}}}}
\def\Set#1{\ensuremath{\left\{\,#1\,\right\}}}
\def\BigSkip{%

\bigskip

}

\def\MedSkip{%

\medskip

}
\def\SmallSkip{%

\smallskip

}
\def\SineWave#1{%
{%
\pgfmathsetmacro{\Xone}{#1}
\pgfmathsetmacro{\Xtwo}{-(#1)}
\draw (0,0) sin(1,\Xone) cos (2,0) sin(3,\Xtwo) cos (4,0);
}
}
\newcommand{\tolstrut}{%
  \vrule height\dimexpr\fontcharht\font`\A+.1ex\relax width 0pt\relax
}

\DeclareRobustCommand{\Overline}[1]{%
  \ensuremath{\overline{\raisebox{0pt}[1.2\height]{#1}}}%
}
\newenvironment{TIKZP}[1][scale=1.0]{%
\adjustbox{valign=t}\bgroup
\begin{tikzpicture}[#1]
}{%
\end{tikzpicture}
\egroup
}
\long\def\Definition#1{%
\begin{block}{Definition}
#1
\end{block}}
\long\def\Postulate#1#2{%
\begin{block}{Postulate}
\begin{description} \item[#1] #2\end{description}
\end{block}}
\long\def\TwoUnequalColumns#1#2#3#4{%
\begin{columns}%
\begin{column}{#1}#3\end{column}%
\begin{column}{#2}#4\end{column}%
\end{columns}%
}
\long\def\ThreeUnequalColumns#1#2#3#4#5#6{%
\begin{columns}%
\begin{column}{#1}#4\end{column}%
\begin{column}{#2}#5\end{column}%
\begin{column}{#3}#6\end{column}%
\end{columns}%
}
\long\def\ThreeColumns#1#2#3{%
\ThreeUnequalColumns{0.33\textwidth}{0.33\textwidth}{0.33\textwidth}{#1}{#2}{#3}%
}
\long\def\TwoColumns#1#2{%
\TwoUnequalColumns{0.5\textwidth}{0.5\textwidth}{#1}{#2}%
}
\long\def\OnlyRemark#1#2{%
\only<#1>{\Remark{#2}}}

\long\def\Remark#1{%
\begin{block}{Remark}
#1
\end{block}}
\long\def\Example#1{%
\begin{example} #1\end{example}}
\def\VV{\textit{vice versa}}
\def\QZero{\ket{0}}
\def\QOne{\ket{1}}
\def\QState#1{\ket{\QName{#1}}}
\def\QName#1{\ensuremath{\psi_{#1}}}
\def\TensProd#1#2{\ensuremath{#1\otimes #2}}
\def\TensSupProd#1#2{\ensuremath{#1^{\otimes #2}}}
\def\PZero{\SQB{1}{0}}
\def\POne{\SQB{0}{1}}
\def\PHad{\ensuremath\frac{1}{\sqrt{2}}\SQB{1}{1}}
\def\RootTwo{\ensuremath{\frac{1}{\sqrt{2}}}}
\def\Intersect#1#2{\ensuremath{#1 \cap #2}}
\def\Union#1#2{\ensuremath{#1 \cup #2}}
\def\Set#1{\ensuremath{\left\{\,#1\,\right\}}}
\def\AllBits#1#2{\ensuremath{#1\in \Set{0,1}^{#2}}}
\def\DQB#1#2#3#4{%
\ensuremath{\begin{pmatrix*}[r] #1 \\ #2 \\ #3 \\ #4\end{pmatrix*}}}
\def\QQB#1#2#3#4#5#6#7#8{%
\ensuremath{\begin{pmatrix*}[r] #1 \\ #2 \\ #3 \\ #4 \\ #5 \\ #6 \\ #7 \\ #8\end{pmatrix*}}}
\def\SQB#1#2{%
\ensuremath{\begin{pmatrix*}[r] #1 \\ #2\end{pmatrix*}}}
\def\TZPointDiam{2pt}
\def\TZPoint#1#2#3{%
\draw[fill=black] (#1) circle (\TZPointDiam) node[#3] {#2} ;}
\def\UnitComplexCircle{%
\draw [<->] (-1.5, 0) -- (1.5, 0)  node[right] {$\Re{}$} ;
   \draw [<->] (0,-1.5) -- (0, 1.5) node[above] {$\Im{}$} ;

   \draw (0, 0) circle (1) ;
   }
   \def\TZPEast{\TZPoint{1,0}{$+1$}{above right}}
   \def\TZPNorth{\TZPoint{0,1}{$i$}{above right}}
   \def\TZPWest{\TZPoint{-1,0}{$-1$}{above left}}
   \def\TZPSouth{\TZPoint{0,-1}{$-i$}{below right}}
\def\TZText#1#2#3{%
  \draw (#1) node [#3] {#2};
}
\def\Qiskit{\texttt{qiskit}}
\def\TwoToThe#1{\ensuremath{2^{#1}}}
\def\Exp#1{\ensuremath{e^{#1}}}
\def\ExpPhase#1{\Exp{i\,#1}}
\def\ExpNegPhase#1{\Exp{-i\,#1}}
\def\PhaseLabel#1{\phase{\fbox{#1}}}
\def\ColHead#1{\begin{center}\bf #1\end{center}}
\def\Complex#1#2{\ComplexGen{#1}{#2}{+}}
\def\ComplexDiff#1#2{\ComplexGen{#1}{#2}{-}}
\def\ImaginaryM#1{\ComplexDiff{\relax}{#1}}
\def\ComplexGen#1#2#3{%
\ensuremath{#1 #3 i\relax #2}}
\def\Polar#1#2{\ensuremath{#1\hbox to 0.1pt{\hss}\Exp{i\relax#2}}}
%%
%% #1 -- lower left x
%% #2 -- lower left y
%% #3 -- upper right x
%% #4 -- upper right y
%% #5 -- num lines
\def\PFilter#1#2#3#4#5{

    \draw (#1,#2) -- (#3,#2) -- (#3,#4) -- (#1,#4) -- cycle;
    \foreach \i in {1,...,#5}
    {
        \pgfmathsetmacro{\PFilterInc}{#1 + (#3-#1) / #5 * \i}
        % \edef\PFilterInc{\pgfmathresult}
        \draw (\PFilterInc, #2) -- (\PFilterInc, #4);
    }
}
%% #1 - x
%% #2 - y
%% #3 - degrees
%% #4 - options
%% #5 - radius
\def\RadiantArrowsR#1#2#3{%
\foreach \p in {0, #1,...,360} {
   \draw[#2] (0,0)  -- (\p:#3) ;
}
}
\def\RadiantArrows#1#2{%
   \RadiantArrowsR{#1}{#2}{1}
}
\newcommand\RotateTheta[1][scale=1.0]{%
\SquareOutline{}
\begin{scope}[#1]
\fill[white] (0.5,0.5) circle (0.5);
\draw (0.5,0.5) circle (0.5);
\draw[->] (1.0,0.5) arc(0:45:0.5);
\draw (0.5,0.5) node {$\theta$};
\end{scope}
}
\def\SquareOutline{%
\path (0,0) rectangle (1,1);
}
\newcommand\Qif[1][scale=1.0]{%
\begin{scope}[#1]
\SquareOutline{}
\clip (0,0) rectangle(1,1);
\RotateAroundCenter{45}{%
\begin{scope}
\draw (0.15,0.15) rectangle(0.85,0.85);
\end{scope}}
\end{scope}
}
\newcommand\LightSource[1][scale=1.0]{%
\begin{scope}[#1]
\SquareOutline{}
\PhotonGun{}
\begin{scope}[shift={(0.75,0.5)}]
    \RadiantArrowsR{22.5}{color=red}{0.25}
\end{scope}
\end{scope}
}
\def\PhotonGun{%
\fill (0,0.375) rectangle (0.5, 0.625);
}
\def\HPolar{%
\begin{scope}[color=red]
\SquareOutline{}
\draw[thick] (0.5,0.4) -- ++(0.35,0)
             ++(-0.35,0.1) -- ++(0.35,0)
             ++(-0.35,0.1) -- ++(0.35,0);
% \fill (0.90,0.5) circle (0.05);
\end{scope}
}
\def\VPolar{%
\begin{scope}[color=red]
\draw[thick] (0.6,0.375) -- ++(0,.250)
             ++(0.1,-.250) -- ++(0,.250)
             ++(0.1,-.250) -- ++(0,.250);
% \fill (0.95,0.5) circle (0.05);
\end{scope}
}
\newcommand\HPolarizedLightSource[1][scale=1.0]{%
\begin{scope}[#1]
\SquareOutline{}
\PhotonGun{}
\HPolar{}
\end{scope}
}
\newcommand\VPolarizedLightSource[1][scale=1.0]{%
\begin{scope}[#1]
\SquareOutline{}
\PhotonGun{}
\VPolar{}
\end{scope}
}
\newcommand\EVBomb[1][scale=1.0]{%
\begin{scope}[#1]
\fill (0.5,0.5) circle(0.5);
\end{scope}%
}
\newcommand\EVBoom[1][scale=1.0]{%
\begin{scope}[#1]
\fill[white] (0.5,0.5) circle(0.5);
\draw[red,thick] (0.5,0.5) circle(0.5);
\clip (0,0) rectangle(1,1);
\begin{scope}
\node[scale=0.35,xshift=28pt,yshift=28pt,starburst,draw,red,fill=orange,line width=1.5pt]
{BOOM!};
\end{scope}
\end{scope}
}
\newcommand\Measurement[1][scale=1.0]{%
\begin{scope}[#1]
\draw[color=white] (0,0) rectangle(1,1);
\begin{scope}
\clip (0.25,0.25) rectangle (1,1);
\draw (0,0) circle(0.9);
\draw (0,0) circle(1.0);
\end{scope}
\draw[thick] (0,0) --(0,1);
\draw[->] (0.25,0.25) -- (0.9,0.9);
\end{scope}
}
\def\TRectangle#1#2#3#4#5#6{%
\begin{scope}[#6]
\draw [fill=#5] (#1,#2) rectangle (#3,#4);
\end{scope}
}
%% #1 -- llx
%% #2 -- lly
%% #3 -- urx
%% #4 -- ury
%% #5 -- angle
\newcommand\Mirror[1][rotate=0]{%
\clip (0,0) rectangle(1,1);
\begin{scope}[#1]
   \SquareOutline{}
   \TRectangle{0}{.4375}{1}{.5625}{Gray}{rotate=0}
\end{scope}
}
\newcommand\BeamSplitter[1][scale=1.0]{%
\begin{scope}[#1]
    \SquareOutline{}
    \TRectangle{0}{.375}{1}{.625}{SkyBlue}{rotate=0}
\end{scope}
}
\def\MZ{%
\LightSource{}
\Shift{2}{0}{\RotateAroundCenter{-45}{\BeamSplitter{}}}
\Shift{2}{-2}{\RotateAroundCenter{-45}{\Mirror{}}}
\Shift{6}{0}{\RotateAroundCenter{-45}{\Mirror{}}}
\Shift{6}{-2}{\RotateAroundCenter{-45}{\BeamSplitter{}}}
\Shift{8}{-2}{\Measurement[color=orange]{}}
\Shift{6}{-4}{\RotateAroundCenter{-90}{\Measurement[color=NavyBlue]{}}}
}
\def\Shift#1#2#3{%
\begin{scope}[shift={(#1,#2)}] #3\end{scope}
}
\def\RotateAroundCenter#1#2{%
\begin{scope}[rotate around={#1:(0.5,0.5)}]
  #2
\end{scope}
}
\def\Vskip#1{\mbox{}\vskip #1\mbox{}}
\def\Hskip#1{\mbox{}\hskip #1\mbox{}}
\author{Ron K.~Cytron}
\institute{Washington University\\Saint Louis, Missouri\\[2em] \textcopyright~2021 All rights reserved by the author}
\setbeamertemplate{footline}[frame number]{}
%\setbeamertemplate{footline}[text line]{%
%  \parbox{\linewidth}{\vspace*{-8pt}\textcopyright~2021\hfill\insertshortauthor\hfill\insertpagenumber}}
%% #1 -- lecture number
%% #2 -- lecture title
%% #3 -- lecture subtitle
%% #4 -- lecture label
\def\SetTitle#1#2#3#4{%
   \lecture[#1]{#2}{lec:#4}%
   \title{#2}%
   \subtitle{#3}%
   \date{\today}%
   \begin{frame}\maketitle\end{frame}%
}
\def\Kaye{\href{https://dl.acm.org/doi/10.5555/1206629}{Kaye}}
\def\MikeIke{\href{https://dl.acm.org/doi/10.5555/1972505}{Nielson~and~Chuang}}

\newcounter{ProtocolDialogStep}
\newenvironment{ProtocolDialog}[3]{%
\def\Incr{\stepcounter{ProtocolDialogStep}}
\def\Ref##1{\Incr{}\ThreeUnequalColumns{#1}{#2}{#3}{\arabic{ProtocolDialogStep}. ##1}{\relax}{\relax}}%
\def\Alice##1{\Incr{}\ThreeUnequalColumns{#1}{#2}{#3}{\relax}{\arabic{ProtocolDialogStep}. ##1}{\relax}}%
\def\All##1##2##3{%
\Incr{}
\ThreeUnequalColumns{#1}{#2}{#3}{\arabic{ProtocolDialogStep}. ##1}{##2}{##3}
}
\def\Bob##1{\Incr{}\ThreeUnequalColumns{#1}{#2}{#3}{\relax}{\relax}{
\arabic{ProtocolDialogStep}. ##1}}%
\setcounter{ProtocolDialogStep}{0}%
\ThreeUnequalColumns{#1}{#2}{#3}{\textbf{Ref}}{\textbf{Alice}}{\textbf{Bob}}
}{%
}
%% #1 -- options
%% #2 -- width
%% #3 -- height
%% #4 -- signals
\newenvironment{GateBox}[4][scale=1.0]{%
\edef\Signals{#4}
\edef\Height{#3}
\edef\Width{#2}
\pgfmathsetmacro{\Max}{\Signals}
\pgfmathsetmacro{\Vsep}{\Height / \Max}
\pgfmathsetmacro{\Wlen}{\Width / 5}
\def\CalcY##1{%
0}
\def\Input##1##2{%
   \pgfmathsetmacro{\Xone}{-\Wlen}
   \pgfmathsetmacro{\Xtwo}{0}
   \pgfmathsetmacro{\MyY}{\Height - 0.5 * \Vsep - ##1*\Vsep}
   \draw[->] (\Xone,\MyY) node[left] {##2} -- (\Xtwo,\MyY);
}
\def\Output##1##2{%
   \pgfmathsetmacro{\Xone}{\Width}
   \pgfmathsetmacro{\Xtwo}{\Width + \Wlen}
   \pgfmathsetmacro{\MyY}{\Height - 0.5 * \Vsep - ##1*\Vsep}
   \draw[->] (\Xone,\MyY) -- (\Xtwo,\MyY) node[right] {##2};
}
\def\BoxLabel##1{%
\node[draw, fit={(0,0) (\Width,\Height)}, inner sep=0pt, label=center:\mbox{##1}] (A) {};
}
\begin{TIKZP}[#1]
   \draw (0,0) rectangle(#2,#3);
}{%
\end{TIKZP}
}
\def\PhiTheta#1#2{%
\mbox{\ensuremath{\textcolor{OrangeRed}{\phi = #1}, \textcolor{NavyBlue}{\theta = #2}}}}
\newenvironment{BlochSphere}[3][scale=1.0]{
    \bgroup%
    \def\BlochSpherePhiTheta{%
    \tdplotsetrotatedcoords{0}{0}{0};
        \path[tdplot_rotated_coords] (1.0,0,0) arc (0:115:1.0) node(phi) {};
        \draw[tdplot_rotated_coords,thick,OrangeRed] (0.3,0,0) arc (0:115:0.3);
        \draw[tdplot_rotated_coords,OrangeRed,thick,dashed] (0,0,0) -- (phi);
        \draw[OrangeRed] (0.4,0.4,0) node {$\phi$};
        
        \tdplotsetrotatedcoords{25}{-90}{0};
        \path[tdplot_rotated_coords] (1,0,0) arc (0:45:1) node(theta) {};
        \draw[tdplot_rotated_coords,thick,NavyBlue] (0.4,0,0) arc (0:45:0.4);
        \draw[tdplot_rotated_coords,NavyBlue,thick] (0,0,0) -- (theta);
        \draw[NavyBlue] (-0.4,0,0.3) node {$\theta$};
        
        
    }
    \def\TZPointDiam{0.7pt}%
    \def\ShowAxes{%
    \draw[dashed, tdplot_rotated_coords, black] (1,0,0) node {{\Large\bf x}} (0,1,0) node {{\Large\bf y}} (0,0,1) node {{\Large\bf z}};
    }%
    \tdplotsetmaincoords{#2}{#3}%
    \begin{tikzpicture}[tdplot_main_coords, #1]
        
        
        \shade[ball color = lightgray, opacity = 0.5] (0,0,0) circle (1cm);
        
        %%
        %% The equator
        %%
        \tdplotsetrotatedcoords{0}{0}{0};
        \draw[dashed, tdplot_rotated_coords, gray] (0,0,0) circle (1);
        
        %%
        %% Arc in the x-z plane
        %%
        \tdplotsetrotatedcoords{90}{90}{90};
        \draw[dashed, tdplot_rotated_coords, gray] (1,0,0) arc (0:180:1);
        
        %%
        %% Arc in the y-z plane
        %%
        \tdplotsetrotatedcoords{0}{90}{90};
        \draw[dashed, tdplot_rotated_coords, gray] (1,0,0) arc (0:180:1);
        
        %\tdplotsetrotatedcoords{90}{90}{90};
       % \draw[dashed, tdplot_rotated_coords, black] (1,0,0) node {{\Large\bf x}} (0,1,0) node {{\Large\bf y}} (0,0,1) node {{\Large\bf z}}; 
        
        \draw[dashed, gray] (0,0,0) -- (-1,0,0);
        \draw[dashed, gray] (0,0,0) -- (0,-1,0);
        \draw[dashed, gray] (0,0,0) -- (0,0,-1);
        
        \draw[-stealth] (0,0,0) -- (1.30,0,0);
        \draw (1.45,0,0) node {$x$};
        \draw[-stealth] (0,0,0) -- (0,1.30,0);
        \draw(0,1.45,0) node {$y$};
        \draw[-stealth] (0,0,0) -- (0,0,1.30);
        \draw (0,0,1.45) node {$z$};
        
}{%
    \end{tikzpicture}
    \tdplotsetmaincoords{0}{0}
    \egroup
}

\long\def\Exercise#1#2{%
\relax}
\def\QED{\begin{flushright}\ensuremath{\square}\end{flushright}}


%%
%% test
%%
\begin{document}
\SetTitle{6}{A single qubit}{Representation and some quantum gates}{06}

\section{Overview}
\begin{frame}{Overview}{What we will cover}
\begin{itemize}
    \item We develop an efficient representation for a single quantum bit (qubit), parameterized by two real values.
    \item We introduce the Bloch sphere: a compelling visualization of a single qubit.
    \item We introduce the Hadamard gate that transforms \ket{0} into a uniform superposition of \ket{0} and \ket{1}.
    \item We introduce the \emph{Pauli} gates and their associated unitary matrices.
    \item We develop a general approach to rotating a qubit around each axis of the Bloch sphere.
\end{itemize}
\end{frame}

\section{Efficient characterization}

\begin{frame}{State space}{Definition}
\Vskip{-3em}\Postulate{State Space}{The state of a quantum system is a unit vector in a complex \href{https://en.wikipedia.org/wiki/Hilbert_space}{Hilbert space}}
\begin{itemize}
    \item We need a Hilbert space, rather than the simpler \href{https://en.wikipedia.org/wiki/Euclidean_space}{Euclidean space}, because we use complex values to represent state.
    \item We have used two dimensions to specify a single qubit, one each for $\alpha$ and $\beta$: $\ket{\psi}=\SQB{\alpha}{\beta}$, where
    $\Prob{\alpha}+\Prob{\beta}=1$.
  
    \item The state of an $n$-qubit system is a unit column vector (ket)  of $\TwoToThe{n}$ complex values.
\end{itemize}
\OnlyRemark{1}{%
We next examine a more concise and visually compelling definition of a qubit.
}
\end{frame}

\begin{frame}{Coordinates for a qubit}{Cartesian and polar}
\Vskip{-5em}\TwoUnequalColumns{0.4\textwidth}{0.6\textwidth}{%
\ColHead{Cartesian}
\begin{itemize}
    \item $\ket{\psi}=\alpha\ket{0} + \beta\ket{1}$ with $\Mag{\alpha}^2+\Mag{\beta}^2=1$.
    \item Letting
    \begin{itemize}
        \item $\alpha=\Complex{a_{0}}{b_{0}}$
        \item $\beta=\Complex{a_{1}}{b_{1}}$
    \end{itemize}
    \item This requires \alert{four real parameters}
\end{itemize}
}{%
\ColHead{Polar}
\begin{itemize}
    \item Write $\alpha = r_{a}\left(\cos\theta_{a}+\NiceI \sin\theta_{a}\right) = \Polar{r_{a}}{\theta_{a}}$
    \item Write $\beta = r_{b}\left(\cos\theta_{b}+\NiceI \sin\theta_{b}\right) = \Polar{r_{b}}{\theta_{b}}$
    \item Then \begin{eqnarray*} \ket{\psi} & = & \Polar{r_{a}}{\theta_{a}}\ket{0} + \Polar{r_{b}}{\theta_{b}}\ket{1} \\
    & = & \visible<1>{\ExpPhase{\theta_{a}}}\left( \alert<4>{r_{a}}\ket{0} + \Polar{\alert<4>{r_{b}}}{
    \only<1-2>{(\theta_{b}-\theta_{a})}
    \only<3->{\alert<4>{\phi}}
    }\ket{1} \right)
    \end{eqnarray*}
    \item<2-> Discard $\ExpPhase{\theta_{a}}$ as a global phase
    \item<3-> Let $\phi=\theta_{b}-\theta_{a}$
    \item<4-> So we need just \alert{three real parameters}
\end{itemize}
}
\end{frame}

\def\Bloch#1#2{%
\ensuremath{%
\cos{\frac{\textcolor{\RCtwo}{#1}}{2}}\ket{0} + \Polar{}{\textcolor{\RCone}{#2}}\sin{\frac{\textcolor{\RCtwo}{#1}}{2}}\ket{1}}}

\section{Bloch sphere}

\begin{frame}{Can we do better?}{Yes!  It takes only two real parameters to denote the state of one qubit}
\begin{itemize}
   \item We start with three real parameters:
   \( \ket{\psi}= r_{a}\left( \ket{0} + r_{b}\ExpPhase{\phi}\ket{1}\right) \)
   \item We have not yet accounted for the unit nature of \ket{\psi}.   
   \item We require $r_{a}^{2}+r_{b}^{2}=1$ which implies \And{r_{a}^{2} \leq 1}{r_{b}^{2} \leq 1}. 
   \item We can show $\Implies{x^{2}\leq 1}{-1 \leq x \leq 1}$.
   \item<2-> We therefore substitute for $r_{a}$ and $r_{b}$ as follows:
   \[
     r_{a} = \cos{\frac{\theta}{2}}, 
     r_{b} = \sin{\frac{\theta}{2}} \]
   The use of $\theta/2$ is arbitrary for now but will work out well later.
   \item<3-> We can thus characterize $\ket{\psi}$ using only \alert{two (angular) real parameters}:
   \[ \ket{\psi} = \Bloch{\theta}{\phi}
   \]
\end{itemize}

    
\end{frame}

\begin{frame}{From polarizing filters to a sphere}{The angular separation between the basis states changes}
\Vskip{-3em}\TwoColumns{%
\begin{itemize}
    \item<1-> With polarizing filters, our bases were horiziontal \visible<2->{and vertical, so they were separated by \Degrees{90}.}
    \item<3-> We imposed a \Degrees{90} angular separation between these bases logically as well.
    \item<4-> We are about to use a \emph{sphere} to capture quantum state.  To use the entire sphere, we locate \ket{0} at the north pole and \ket{1} at the south pole.
    \item<5-> In that depiction, the basis states are separated by \Degrees{180}.
\end{itemize}
}{%
\Vskip{-3em}\begin{center}
\begin{TIKZP}
\visible<3->{
\draw[->,ultra thick] (0,0) -- (1,0) node[right] {\ \ket{0}};
\draw[->,ultra thick] (0,0) -- (0,1) node[above] {\ket{1}};
}
\visible<1->{\begin{scope}[draw=\RCtwo,rotate=90]\PFilter{-1}{-1.2}{1}{1.2}{20}\end{scope}}
\visible<2->{\begin{scope}[draw=\RCthree]\PFilter{-1}{-1.2}{1}{1.2}{20}\end{scope}}
\end{TIKZP}
\end{center}
\visible<4->{%
\begin{center}
    \begin{BlochSphere}[scale=1]{60}{115}
    \BlochSpherePhiTheta{}
     \TZPoint{0,0,1}{\ket{0}}{right}
     \TZPoint{0,0,-1}{\ket{1}}{below}
    \end{BlochSphere}
\end{center}}
}
\end{frame}

\begin{frame}{Our two-parameter characterization}{It meshes well with our studies so far}
\[ \Bloch{\theta}{\phi} \]
\begin{itemize}
    \item The parameter \textcolor{\RCtwo}{$\theta$} trades wave amplitude between \ket{0} and \ket{1}.
    \item Only the relative phase between \ket{0} and \ket{1} matters for interactions.  
    \item This formulation ascribes the phase difference to \ket{1}, using a walk of \textcolor{\RCone}{$\phi$} radians counterclockwise around the complex unit circle, starting at 0~radians.
    \item This formulation of state makes clear that while only two real-valued parameters are needed, the number of states is \href{https://en.wikipedia.org/wiki/Uncountable_set}{uncountably infinite}.
    
\end{itemize}
    
\end{frame}

\begin{frame}{Recall the beam splitter}{We characterize two of its states using our new formula}
\[ \Bloch{\theta}{\phi} \]
    \begin{itemize}
        \item The state of the photon was initially $\Bloch{0}{0} = \SQB{1}{0}$
        \item After one beam splitter we obtained \[ \Bloch{\pi/2}{\pi/2} = \frac{1}{\sqrt{2}} \SQB{1}{i} \]
    \end{itemize} 
\end{frame}

\begin{frame}{The Bloch sphere}{Our two angular parameters determine a point on the unit sphere}
\TwoColumns{%
\[ \ket{\psi} = \Bloch{\theta}{\phi} \]
\begin{itemize}

\item \ket{\psi} is located 
\begin{itemize}
    \item  \textcolor{\RCone}{$\phi$} radians off the $x$~axis and 
    \item \textcolor{\RCtwo}{$\theta$} radians off the $z$~axis.
\end{itemize}
\item We locate states \emph{on} the unit sphere, but the distance from the center of the sphere has no meaning here.  

A more general interpretation of the sphere ascribes meaning to interior points.
\end{itemize}
}{%
\Vskip{-3em}\begin{BlochSphere}[scale=2.5]{60}{115}
    
    \BlochSpherePhiTheta{}
    \TZPoint{theta}{\ket{\psi}}{above right}
\end{BlochSphere}
}
\end{frame}

\begin{frame}{Correspondence between the formula and the sphere}{We see why the formula cuts the angle in half}
\TwoUnequalColumns{0.5\textwidth}{0.5\textwidth}{%
\begin{itemize}
  \item<2-> We locate \ket{0} at \PhiTheta{0}{0} using \Bloch{0}{0}.  The choice of \textcolor{\RCone}{$\phi$} doesn't matter when $\textcolor{\RCtwo}{\theta}=0$.
  \item<3-> We locate \ket{1} at $\PhiTheta{0}{\pi}$ using $\Bloch{\pi}{0}$.
  \item<4-> From the picture, \textcolor{\RCone}{$\phi$} doesn't matter when $\textcolor{\RCtwo}{\theta}=\pi$.  How can this be?
\end{itemize}
}{%
\only<1-4>{%
\[ \ket{\psi} = \Bloch{\theta}{\phi} \]}
\only<5>{%
\[ \ket{\psi} =  \ExpPhase{\textcolor{\RCone}{\phi}} \ket{1} \]
}%
\only<1-4>{%
\begin{BlochSphere}[scale=1.5]{60}{115}
    
    \BlochSpherePhiTheta{}
    \visible<2->{\TZPoint{0,0,1}{$\ket{0}$}{right}}
    \visible<3->{\TZPoint{0,0,-1}{$\ket{1}$}{right}}
    
\end{BlochSphere}}%
\only<5>{%
\begin{itemize}
\item With $0$ as the amplitude on \ket{0}, any phase on \ket{1} with $\textcolor{\RCone}{\phi}\not=0$ can be canceled with an appropriate global phase of 
  \ExpPhase{(-\textcolor{\RCone}{\phi})} on \ket{\psi}.
\item Global phases can be ignored as they do not affect measurements or interactions.
\end{itemize}
}
}
\end{frame}

\begin{frame}{The computational basis}{Measurement outcomes are related to proximity to the north or south pole}
\TwoUnequalColumns{0.6\textwidth}{0.4\textwidth}{%
\begin{itemize}
  \item<1-> The north and south poles, \ket{0} and \ket{1}, respectively, are our \emph{standard} or \emph{computational basis}.
   All other points on the sphere are superpositions of that basis.
   \item<2-> The probability of measuring \ket{0} or \ket{1} is related to how close the state is to the north or south pole, respectively, which is represented by $\textcolor{\RCtwo}{\theta}$.
   \item<3-> The probability of measuring \ket{0} is $\cos^{2}(\textcolor{\RCtwo}{\theta}/2)$.  At the equator, we are equally likely to measure \ket{0} or \ket{1}.
\end{itemize}
}{%
\[ \ket{\psi} = \Bloch{\theta}{\phi} \]
\begin{BlochSphere}[scale=1.5]{60}{115}
    
    \BlochSpherePhiTheta{}
    \visible<2->{\TZPoint{0,0,1}{$\ket{0}$}{right}}
    \visible<3->{\TZPoint{0,0,-1}{$\ket{1}$}{right}}
    
\end{BlochSphere}
}
\end{frame}

{%
\def\Sket#1{%
\ket{\hbox to 4ex{\hfil\ensuremath{#1}\hfil}}}%
\def\R#1{\mbox{\textcolor{\RCone}{#1}}}%
\def\B#1{\mbox{\textcolor{\RCtwo}{#1}}}%
\begin{frame}{Navigating the Bloch sphere}{Other points of interest}


\Vskip{-4em}\TwoUnequalColumns{0.6\textwidth}{0.4\textwidth}{%
\begin{center}
    \begin{tabular}{ccc}
        State & \textcolor{\RCone}{$\phi$} & \textcolor{\RCtwo}{$\theta$} \\
        \visible<1->{\Sket{0} & \R{$0$} & \B{$0$} \\}%
        \visible<1->{\Sket{1} & \R{$0$} & \B{$\pi$} \\}%
        \visible<2->{\Sket{+x} & \R{$0$} & \B{$\pi/2$} \\}%
        \visible<3->{\Sket{-x} & \R{$0$} & \B{$-\pi/2$} \\}%
        \only<3-4>{\Sket{-x} & \R{$\pi$} & \B{$\pi/2$} \\}%
        \visible<5->{\Sket{+y} & \R{$\pi/2$} & \B{$\pi/2$} \\}%
        \visible<6->{\Sket{-y} & \R{$-\pi/2$} & \B{$\pi/2$} \\}%
    \end{tabular}
\end{center}
\only<4>{\Vskip{-5em}States \ket{+x} and \ket{-x} are also called \ket{+} and \ket{-}, respectively.  They are more familiar as:
\begin{eqnarray*} \ket{+} & = & \frac{1}{\sqrt{2}} \SQB{1}{1} \\
\ket{-} & = & \frac{1}{\sqrt{2}} \SQB{1}{-1}
\end{eqnarray*}
}%

}{%
\[ \ket{\psi} = \Bloch{\theta}{\phi} \]
\begin{BlochSphere}[scale=1.5]{60}{115}
    
    \BlochSpherePhiTheta{}
    \visible<1>{\TZPoint{0,0,1}{$\ket{0}$}{right}}
    \visible<1>{\TZPoint{0,0,-1}{$\ket{1}$}{right}}
    \visible<2-3>{\TZPoint{1,0,0}{$\ket{+x}$}{right}}
    \visible<3>{\TZPoint{-1,0,0}{$\ket{-x}$}{right}}
    \visible<5-6>{\TZPoint{0,1,0}{$\ket{+y}$}{above}}
    \visible<6>{\TZPoint{0,-1,0}{$\ket{-y}$}{right}}
\end{BlochSphere}
}
\Exercise{y}{What are +y and -y, expressed in the standard basis?}
\end{frame}
}

\begin{frame}{Switching between bases}{We consider here the standard and Hadamard bases}
\TwoUnequalColumns{0.7\textwidth}{0.3\textwidth}{%
\begin{itemize}
    \item Typically, our computations begin and end (with measurements) in the standard basis.
    \item We often create a uniform superposition of the standard basis states, \ket{+} or \ket{-}, which together form the \href{https://en.wikipedia.org/wiki/Quantum_logic_gate\#Hadamard_gate}{Hadamard} basis in quantum computing.
    \item<1-> How?  We need a transformation (unitary matrix) \Hadamard{} that sends $\ket{0}\mapsto \ket{+}$ and $\ket{1}\mapsto \ket{-}$.
    \item<3-> From linear algebra, the columns of \Hadamard{} are the transformed versions of \textcolor<4->{\RCtwo}{\ket{0}} and \textcolor<5->{\RCthree}{\ket{1}}.
    \item<6-> Because it is symmetric and its entries are real, \mbox{\Hadamard{} = \Conj{\Hadamard{}}}  $\rightarrow$ \Hadamard{} is also its own inverse.
    
\end{itemize}
}{%
\visible<2->{%
\begin{eqnarray*}
    \SQB{1}{0} & \mapsto & \textcolor<3->{\RCtwo}{\frac{1}{\sqrt{2}}\SQB{1}{1}} \\
    \SQB{0}{1} & \mapsto & \textcolor<3->{\RCthree}{\frac{1}{\sqrt{2}}\SQB{1}{-1}}
    \end{eqnarray*}}
\visible<4->{%
\[
\Hadamard{} = \HMatrix{}
\]
\begin{TIKZP}[overlay]
  \draw[white] (0,0) circle(1pt);  % anchor
  \draw<4-5>[\RCtwo,thick] (2.25,0.7) ellipse(0.3 and 0.6);
  \draw<5>[\RCthree,thick] (3, 0.7) ellipse (0.3 and 0.6);
\end{TIKZP}
}
}
\end{frame}

\begin{frame}{Physics and the Bloch sphere}{Models electron spin in space}
\Vskip{-4em}\TwoUnequalColumns{0.8\textwidth}{0.2\textwidth}{%
\begin{itemize}
\item The labeled axes on the Bloch sphere are often used in physics to model spin in each of the three physical dimensions:
\begin{description}
\item[z axis] spin up or down
\item[x axis] spin in or out
\item[y axis] spin left or right
\end{description}
\item In quantum computing we almost always use the standard basis of \ket{0} and \ket{1} (the z axis).
\item The named axes are mutually orthogonal. \Exercise{Orthogonal Axes}{Prove orthogonality of the x, y, and z axes}
\item While we can characterize two of the dimensions using only real values (the z--x plane, where $\textcolor{\RCone}{\phi}=0$), we require complex values to admit a third mutually orthogonal dimension (y).

\end{itemize}
}{%
\begin{center}
    \begin{BlochSphere}[scale=1]{60}{115}
    
    \BlochSpherePhiTheta{}
    \end{BlochSphere}
\end{center}
}

\end{frame}

\begin{frame}{What other bases are there?}{Quite a few}
\Definition{Two points on the Bloch sphere are \emph{antipodal} if they determine a line that passes through the center of the sphere.}
\TwoColumns{%
\begin{itemize}
    \item Every pair of antipodal points determines an orthogonal basis for the quantum state of a single qubit. \Exercise{Antipodal points determine a basis}{Prove that antipodal points on the Bloch sphere are orthogonal.}
    \item Using linear algebra, we can express any state \ket{\psi} in any basis. \LinkArrow{https://www.youtube.com/watch?v=P2LTAUO1TdA}
\end{itemize}
}{%
\begin{center}
\begin{BlochSphere}[scale=1.5]{60}{115}
    \BlochSpherePhiTheta{}
\end{BlochSphere}
\end{center}
}
\end{frame}

%% \begin{frame}{Antipodal points are orthogonal}{Leave as exercise?}
%% \ToDo{include or not?}
%% \end{frame}

\section{The Pauli matrices}

\begin{frame}{Overview of the Pauli matrices}{We use these as gates in quantum computation}

\begin{itemize}
    \item There are three \href{https://en.wikipedia.org/wiki/Pauli_matrices}{Pauli matrices} that are named for the three axes of the Bloch sphere.
    \item Each matrix has the property of rotating a state by $\pi$ radians about its associated axis.
    \item Each gate is its own inverse.
    \item We will develop a universality argument:  any quantum gate can be realized using a linear combination of the Pauli matrices.
    \item The gates we study here are quantum gates, used in the middle of a quantum circuit.  We can also think of these gates as \emph{measurement operators} but we study that later.
\end{itemize}
    
\end{frame}

\begin{frame}{Pauli \PauliZ{}}{A gate that changes the sign of \ket{1}}

\TwoUnequalColumns{0.6\textwidth}{0.4\textwidth}{%
\[
\PauliZ{} = \ZMatrix{}
\]
\begin{itemize}
    \item<1-> This matrix transforms a state as follows: \[ \ApplyGate{\PauliZ}{\SQB{\alpha}{\beta}} =  \SQB{\alpha}{-\beta} \]
    \item<2-> However, two states are unaffected by \PauliZ{}:
    \begin{itemize}
       \item $\ApplyGate{\PauliZ}{\left(\ket{0}=\SQB{1}{0}\right)} = \ket{0}$
       \item $\ApplyGate{\PauliZ}{\left(\ket{1}=\SQB{0}{1}\right)} = -\ket{1}\equiv \ket{1}$
    \end{itemize}
\end{itemize}
}{%
\Vskip{-4em}\begin{center}
    \begin{BlochSphere}[scale=1]{60}{115}
    \BlochSpherePhiTheta{}
     \TZPoint{0,0,1}{\ket{0}}{right}
     \TZPoint{0,0,-1}{\ket{1}}{below}
    \end{BlochSphere}
\end{center}
\only<2>{%
\Vskip{-2.5em}Such states are called eigenstates (or \href{https://en.wikipedia.org/wiki/Eigenvalues_and_eigenvectors}{eigenvectors}) of \PauliZ{}.  A matrix transforms an eigenstate into a multiple $\lambda$ of itself;  $\lambda$ is called an eigenvalue of the matrix.  Here, $\PauliZ{}\ket{0} = +1\ket{0}$ and $\PauliZ{}\ket{1}=-1\ket{1}$.
}
\only<3>{%
\Vskip{-2em}\PauliZ{} affects all other states, off the z axis, for example
\begin{eqnarray*}
\ApplyGate{\PauliZ}{\ket{+}} & = & \ket{-} \\
\ApplyGate{\PauliZ}{\ket{-}}  & = & \ket{+}
\end{eqnarray*}
rotations of $\pi$ radians around the z~axis.
}}
\end{frame}

\begin{frame}{Pauli \PauliX{}}{Also called the \NamedGate{NOT} gate}

\TwoUnequalColumns{0.6\textwidth}{0.4\textwidth}{%
\[
\PauliX{} = \XMatrix{}
\]
\begin{itemize}
    \item<1-> This transformation exchanges the amplitudes: \[ \ApplyGate{\PauliX}{\SQB{\alpha}{\beta}} =  \SQB{\beta}{\alpha} \]
    \item<2-> The two eigenstates are \ket{+x} and \ket{-x}. \Exercise{unaffected points}{Show that X doesn't affect +x or -x.}
    \item<3-> This gate maps $\ket{0}\mapsto\ket{1}$ and \VV{}.
\end{itemize}
}{%
\begin{center}
    \begin{BlochSphere}[scale=1]{60}{115}
    \BlochSpherePhiTheta{}
    \TZPoint{1,0,0}{\ket{x}}{left}
    \TZPoint{-1,0,0}{\ket{-x}}{right}
    \end{BlochSphere}
\end{center}
\PauliX{} generally
rotates a state  by $\pi$ radians around the x~axis.
}
\end{frame}

\begin{frame}{Pauli \PauliY{}}{Exchanges amplitudes and imposes a sign change}

\TwoUnequalColumns{0.6\textwidth}{0.4\textwidth}{%
\[
\PauliY{} = \YMatrix{}
\]
\begin{itemize}
    \item<1-> This transformation exchanges the amplitudes and imposes a relative sign change: \[ \ApplyGate{\PauliY}{\SQB{\alpha}{\beta}} =  \SQB{-\NiceI\beta}{i\alpha} = -\NiceI\SQB{\beta}{-\alpha} \equiv \SQB{\beta}{-\alpha} \]
    \item<2-> The two eigenstates are \ket{+y} and \ket{-y}. \Exercise{unaffected points}{Define +y and -y in the standard basis and show that Paul Y doesn't affect +y or -y.}
    \item<3-> This gate maps $\ket{0}\mapsto\ket{1}$ and \VV{}.
\end{itemize}
}{%
\begin{center}
    \begin{BlochSphere}[scale=1]{60}{115}
    \BlochSpherePhiTheta{}
    \TZPoint{0,1,0}{\ket{y}}{above}
    \TZPoint{0,-1,0}{\ket{-y}}{left}
    \end{BlochSphere}
\end{center}
\PauliY{} generally
rotates a state  by $\pi$ radians around the y~axis.  
\visible<3>{\MedSkip{}This gate also maps $\ket{+} \mapsto \ket{-}$ and \VV{}.}
}
\end{frame}

\begin{frame}{Properties of Pauli gates}{These should be apparent from our study so far}
\begin{center}
$\PauliX{}=\XMatrix{}$
$\PauliY{}=\YMatrix{}$
$\PauliZ{}=\ZMatrix{}$
\end{center}
\begin{itemize}
    \item Each matrix is \href{https://en.wikipedia.org/wiki/Involutory_matrix}{involutory} (each is its own inverse).
    \item These relations can be verified by computing the appropriate products:
    \begin{center}
    { \setlength{\tabcolsep}{3pt}%
    \begin{tabular}{rcrl@{\hspace{3em}}rcrl@{\hspace{3em}}rcrl} \setlength{\tabcolsep}{1pt}
        \PauliX{}\PauliY{} & = & \NiceI &\PauliZ{} & \PauliZ{}\PauliX{} & = & \NiceI &\PauliY{} & \PauliY{}\PauliZ{} & = & \NiceI &\PauliX{} \\
        \PauliY{}\PauliX{} & = & $-\NiceI$ &\PauliZ{} &  \PauliX{}\PauliZ{} & = & $-\NiceI$ &\PauliY{} & \PauliZ{}\PauliY{} & = & $-\NiceI$ &\PauliX{}
    \end{tabular}}
    \end{center}
    \item Each matrix has \href{https://en.wikipedia.org/wiki/Eigenvalues_and_eigenvectors}{eigenvalues} $\pm 1$.
    \item Each Pauli matrix rotates all states by $\pi$ radians about its associated axis, except the two states at the intercept of its axis with the unit Bloch sphere.  The eigenstates are therefore those intercepts.
\end{itemize}
\end{frame}

%%
%% This material might move to module 12
%%
\begin{frame}{Eigenvalues of the Pauli gates}{And their corresponding eigenvectors}
\only<1-5>{%
\begin{itemize}
    \item<1-> In quantum computing, a \href{https://en.wikipedia.org/wiki/Measurement_in_quantum_mechanics}{measurement of a quantum system} is modeled by an operator $M$ called an \href{https://en.wikipedia.org/wiki/Observable}{observable}.  We can generally view the measurement as yielding $+1$ or $-1$.
    \item<2-> After observation, the system is in a state $\psi$ such that subsequent applications of $\psi$ certainly yield the same measurement.
    \item<3-> If a system has been measured by an operator $M$ and it entered some state $\psi$, then we must have \[ M\psi = \lambda\psi \] where $\lambda=\pm 1$ is the measured result.
    \item<4-> The above equation is the \href{https://en.wikipedia.org/wiki/Eigenvalues_and_eigenvectors\#Formal_definition}{definition} of an eigenvalue~$\lambda$ and its eigenvector~$\psi$.
    \item<5-> For qubits, measurement operators yield $\pm 1$ and there will be a distinct eigenvector for each of those eigenvalues.
\end{itemize}}%
\only<6->{%
\TwoColumns{%
\only<6-12>{%
\begin{itemize}
    \item<6-> Each Pauli matrix is an observable.
    \item<7-> The \PauliZ{} matrix corresponds to the \href{https://en.wikipedia.org/wiki/Qubit\#Standard_representation}{computational basis}, because its eigenvectors are \visible<8->{\alert<8>{\ket{0}}} \visible<9->{and \alert<9>{\ket{1}}.}
    \item<10-> When an \alert<10>{arbitrary state~$\psi$}\visible<11->{ is measured as \ket{0} or \ket{1}, subsequent applications of \PauliZ{} produce the same eigenvalue and eigenvector.}  
    \item<12-> If $\psi$ happens to be \ket{0} or \ket{1}, then re-measurement must yield same state.
\end{itemize}}%
\only<13->{%
\begin{itemize}
    \item<13-> Consider the observable \PauliX{} with eigenvectors \visible<14->{\alert<14>{\ket{+}}} \visible<15->{and \alert<15>{\ket{-}}.}
\end{itemize}
}
}{%
\only<8-12>{%
\begin{alignat*}{5}
    \ZMatrix{} & \alert<8>{\PZero{}} & \mbox{ = } & +1  & \alert<8>{\PZero{}}  \\
    \visible<9->{\ZMatrix{} & \alert<9>{\POne{}} & \mbox{ = } & -1  & \alert<9>{\POne{}} }
\end{alignat*}}%
\only<13->{%
\begin{alignat*}{5}
    \XMatrix{} & \alert<14>{\PHad{}} & \mbox{ = } & +1  & \alert<14>{\PHad{}}  \\
    \visible<15->{\XMatrix{} & \alert<15>{\PHadm{}} & \mbox{ = } & -1  & \alert<15>{\PHadm{}} }
\end{alignat*}}
\begin{center}
\begin{TIKZP}
\visible<8->{
\draw[->,ultra thick] (0,0) -- (1,0) node[right] {\ \alert<8>{\ket{0}}};
\draw[->,ultra thick] (0,0) -- (0,1) node[above] {\alert<9>{\ket{1}}};
\visible<10>{\alert{\draw[->,thick] (0,0) -- (30:1) node[above] {$\psi$};}}
\visible<14>{\alert{\draw[->,thick] (0,0) -- (45:1) node[above] {\ket{+}};}}
\visible<15>{\alert{\draw[->,thick] (0,0) -- (-45:1) node[right] {\ket{-}};}}
}
\visible<8>{\begin{scope}[draw=\RCtwo,rotate=90]\PFilter{-1}{-1.2}{1}{1.2}{20}\end{scope}}
\visible<9>{\begin{scope}[draw=\RCthree]\PFilter{-1}{-1.2}{1}{1.2}{20}\end{scope}}
\visible<14>{\begin{scope}[draw=\RCtwo,scale=0.7,rotate=-45]\PFilter{-1}{-1.2}{1}{1.2}{20}\end{scope}}
\visible<15>{\begin{scope}[draw=\RCthree,scale=0.7,rotate=45]\PFilter{-1}{-1.2}{1}{1.2}{20}\end{scope}}
\end{TIKZP}
\end{center}
}
}

\end{frame}
%%
%% End of stuff that might move to module 12
%%
{%
\def\Redify#1{\textcolor{RedOrange}{#1}}%
\def\Bluify#1{\textcolor{\RCtwo}{#1}}%
\def\Term#1{\frac{i^{#1}\theta^{#1}A^{#1}}{#1!}}%
\def\TermA#1{\frac{i\theta^{#1}A}{#1!}}%
\def\TermI#1{\frac{\theta^{#1}\Identity{}}{#1!}}%
\def\TermF#1{\frac{\theta^{#1}}{#1!}}%
\begin{frame}{A useful theorem for expressing rotations}{A matrix extension of \href{https://en.wikipedia.org/wiki/Euler\%27s_formula}{Euler's formula}}


\Vskip{-3em}
\only<1-6>{%
\[  \alert<2>{A^{2}=\Identity{}} \rightarrow \ExpPhase{\theta A} = \Complex{\cos(\theta)\ \Identity{}}{\sin(\theta)\ A} \]}
\TwoColumns{%
Recall the following \href{https://en.wikipedia.org/wiki/Taylor_series}{Taylor series}:
\begin{align*}
\Exp{x} = & 1 + \frac{x}{1!} + \frac{x^2}{2!} + \frac{x^3}{3!} + \cdots \\
\Redify{\cos(x)} = & \Redify{1 - \frac{x^2}{2!} + \frac{x^4}{4!} - \frac{x^6}{6!} +  \cdots} \\ 
\Bluify{\sin(x)} = & \Bluify{\frac{x}{1!} - \frac{x^3}{3!} + \frac{x^5}{5!} - \frac{x^7}{7!} + \cdots}
\only<2>{%
\\ i^{2} = & -1 \\
   i^{3} = & - i \\
   i^{4} = & \ 1
}
\end{align*}
}{%
\Vskip{-4.3em}\begin{align*}
\ExpPhase{\theta A}
\only<1-2>{%
= {} &
\Identity{} + \Term{1} + \alert<2>{\Term{2}} + \Term{3} \\[0.5em] & + \alert<2>{\Term{4}} + \Term{5} + \cdots \\[1em]
}%
\only<2-3>{%
  = {} & 1\,\Identity{} + 
  \TermA{1} - 
  \alert<2>{\TermI{2}} - 
  \TermA{3} \\[0.5em] & + 
  \alert<2>{\TermI{4}} + 
  \TermA{5} \\[1em]}%
\only<4-5>{%
  = {} & \Redify{1\,\Identity{}} + 
  \Bluify{\TermA{1}} - 
  \Redify{\TermI{2}} - 
  \Bluify{\TermA{3}} \\[0.5em] & + 
  \Redify{\TermI{4}} + 
  \Bluify{\TermA{5}}\\[1em]}%
\only<5-6>{%
  = {} & \Identity{}\left[ \Redify{1} - 
   \Redify{\TermF{2}}  
  + \Redify{\TermF{4}} - \cdots\right]
  \\[0.5em] &
  + i\,\left[
  \Bluify{\TermF{1}}
  - \Bluify{\TermF{3}} 
   + 
  \Bluify{\TermF{5}} - \cdots\right]\,A\\[1em]
  }%
\only<6>{%
= {} & \Identity{}\left[ \Redify{\cos(\theta)} \right]
  \\[0.5em] &
  + i\,\left[
  \Bluify{\sin(\theta)}\right]\,A \QEDsym{}
}
\end{align*}
}
\end{frame}
}


\begin{frame}{Arbitrary rotations}{The Pauli matrices rotate by $\pi$ radians;  how can we rotate by $\theta$ radians?}
We can use the extended Euler formula to express rotation about the Bloch x, y, or z axes:
\begin{align*}
R_{x}(\theta) = & \ExpNegPhase{\left(\frac{\theta}{2}\right)\,\PauliX{}}\\[0.7em]
R_{y}(\theta) = & \ExpNegPhase{\left(\frac{\theta}{2}\right)\,\PauliY{}}\\[0.7em]
R_{z}(\theta) = & \ExpNegPhase{\left(\frac{\theta}{2}\right)\,\PauliZ{}}
\end{align*}
where the matrices in the exponents are the Pauli matrices.

When $\theta=\pi$, $\cos(-\pi/2)=0$ and $\sin(-\pi/2)=-1$ $\rightarrow$ $R_a(\pi) = -i\ A \equiv A$, matching the Pauli gates in terms of their rotation by $\pi$ radians about their respective axes.
\end{frame}

\def\EulerP#1{%
\ComplexDiff{\cos\left(\frac{\theta}{2}\right)\Identity}{\sin\left(\frac{\theta}{2}\right)#1}}

\begin{frame}{Matrix formulation of rotating gates}{We use the matrix form of Euler's formula}

\begin{align*}
R_{x}(\theta) = & \ExpNegPhase{\left(\frac{\theta}{2}\right)\,\PauliX{}}%
\visible<2-3>{= \EulerP{\PauliX{}}}%
\visible<3>{= \SQBG{\relax}{\cos\left(\frac{\theta}{2}\right)}{\ImaginaryM{\sin\left(\frac{\theta}{2}\right)}}{\ImaginaryM{\sin\left(\frac{\theta}{2}\right)}}{\cos\left(\frac{\theta}{2}\right)}}\\[0.7em]
%
%
\alert<4>{R_{y}(\theta)} = & \ExpNegPhase{\left(\frac{\theta}{2}\right)\,\PauliY{}}%
\visible<2-3>{= \EulerP{\PauliY{}}}
\visible<3-4>{= \alert<4>{\SQBG{\relax}{\cos\left(\frac{\theta}{2}\right)}{-\sin\left(\frac{\theta}{2}\right)}{\sin\left(\frac{\theta}{2}\right)}{\cos\left(\frac{\theta}{2}\right)}}}\\[0.7em]
%
%
R_{z}(\theta) = & \ExpNegPhase{\left(\frac{\theta}{2}\right)\,\PauliZ{}}%
\visible<2-3>{= \EulerP{\PauliZ{}}}%
\visible<3>{= \SQBG{\relax}{\ExpNegPhase{\frac{\theta}{2}}}{0}{0}{\ExpPhase{\frac{\theta}{2}}}
\equiv \SQBG{\relax}{1}{0}{0}{\ExpPhase{\theta}}}
\end{align*}
\visible<4>{\alert<4>{Given a state with real amplitudes ($\phi=0$ on the Bloch sphere), this rotation shifts probability between \ket{0} and \ket{1} while keeping the amplitudes real.}}
\end{frame}

\begin{frame}{Interior points of the Bloch sphere}{Representing mixed states}
\Vskip{-2.2em}We now draw a distinction between the following two scenarios.  \visible<2->{On the left, we have an unpolarized photon.}  \visible<3->{On the right, we have a photon that is randomly polarized horizontally or vertically.}
\Vskip{-3em}\TwoColumns{%
\begin{center}
\begin{TIKZP}[scale=1.0]
\visible<1->{\RadiantArrows{22.5}{->,color=BrickRed,thick}}
\end{TIKZP}\end{center}
\visible<2->{\ColorOne{\[\QState{} = \QPlus{}\]}}
}{%
\begin{center}
\begin{TIKZP}[scale=0.64]
\visible<1->{
\draw[->,ultra thick] (0,0) -- (1,0) node[right] {\ \ket{0}};
\draw[->,ultra thick] (0,0) -- (0,1) node[above] {\ket{1}};
}
\visible<1->{\begin{scope}[draw=\RCtwo,rotate=90]\PFilter{-1}{-1.2}{1}{1.2}{20}\end{scope}}
\visible<1->{\begin{scope}[draw=\RCthree]\PFilter{-1}{-1.2}{1}{1.2}{20}\end{scope}}
\end{TIKZP}\end{center}
\visible<3->{\Vskip{-3em}\ColorTwo{\begin{align*}
    \QState{} &= \QZero{} &\mbox{ with probability 0.5} \\
    &= \QOne{} &\mbox{ with probability 0.5}
\end{align*}}}
}
\visible<4->{\Vskip{-2em}The \ColorOne{left state} is \href{https://en.wikipedia.org/wiki/Quantum_state\#Pure_states}{pure}, and is found on the surface of the Bloch sphere as \ket{+x}.   The \ColorTwo{right state} is \href{https://en.wikipedia.org/wiki/Quantum_state\#Mixed_states}{mixed}. Recalling \QZero{} and \QOne{} are the north and south poles, respectively, the \ColorTwo{right state} is the vector sum of~0.5\QZero{} and~0.5\QOne{}, which places it statistically at the center of the sphere.}
    
\end{frame}

\begin{frame}{Summary}

\begin{itemize}
    \item Two real parameters suffice for specifying the state of a single quantum bit.
    \item The Bloch sphere is a convenient visualization of a single qubit's state.
    \item We have studied four gates of interest so far:
    \begin{itemize}
        \item The Hadamard gate creates a uniform superposition of \ket{0} and \ket{1} given the input \ket{0}.
        \item The \PauliX{}, \PauliY{}, and \PauliZ{} gates rotate state by $\pi$ radians about their respective axes.
    \end{itemize}
    \item There are matrices that describe arbitrary rotations about the Bloch axes. 
    \item Interior points of the Bloch sphere represent mixed states.
\end{itemize}
\end{frame}

\end{document}
